Пусть $K$ -- произвольное факториальное кольцо. $F = Quot(K)$.

Из факториальности $F$ следует, что для любых элементов $a_1, \ldots, a_k \in F$ существует единственный с точностью до ассоциированности наибольший общий делитель. 

Определим $c(f)$ многочлена $f(x) \in F[x]$ как НОД его коэффициентов. Заметим, что $c(f)$ определено с точностью до ассоциированности. Будем называть $f(x)$ примитивным многочленом, если $c(f) = 1$.\\

Заметим, что если многочлен неразложим, то он примитивен.

\Statement Любой неразложимый элемент в $K$ прост в $K[x]$.

\Proof
Пусть $p$ -- неразложимый элемент в $K$. Рассмотрим $K[x]/(p) \cong (K/(p))[x]$ и заметим, что $K/(p)$ -- область целостности в силу факториальности кольца $K$. Но тогда и $K[x]/(p) \cong (K/(p))[x]$ -- область целостности, поэтому $p$ прост в $K[x]$.
\EndProof

\textbf{Следствие:} Пусть $p \in K[x]$ -- неразложимый элемент, $deg \ p = 0$. Тогда $p$ прост в $K[x]$.

\Statement Пусть $f$ -- примитивный многочлен. Тогда $f$ неприводим над $K[x] \Leftrightarrow f$ неприводим над $F[x]$.

\Proof
\begin{itemize}
    \item $\Leftarrow$ Пусть $f$ приводим над $K[x]$, то есть $f = gh$ для некоторых $g,h \in K[x] \backslash (\{0\} \cup K[x]^*)$, причем $g, h \not\sim 1$. Тогда $deg \ h, deg \ g \geq 1$. Но тогда $g, h \not\in F[x]^*$, что противоречит неприводимости многочлена $f$ над $F[x]$.
    
    \item  $\Rightarrow$ Пусть $f$ приводим над $F[x]$, то есть $f = gh$ для некоторых  $g,h \in F[x] \backslash (\{0\} \cup F[x]^*)$, $deg \ g, deg \ h \geq 1$. Тогда f = $\frac{A}{B}g_2, h_2$, где $g_2, h_2 \in K[x]$ -- примитивные многочлены положительной степени,  $\frac{A}{B} \in F$. Но тогда $Bf = Ag_2h_2$, и в силу примитивности $f \sim g_2h_2$, что противоречит неприводимости многочлена $f$ над $K[x]$.
\end{itemize}
\EndProof

\textbf{Замечание} Аналогично этому доказательству можно получить, что у любого ненулевого многочлена из $K[x]$ существует разложение на неприводимые над $K[x]$ многочлены, пользуясь факториальностью $F[x]$.\\

\textbf{Следствие:} Пусть $f \in K[x]$ -- неприводимый многочлен, $deg \ f \geq 1$. Тогда $f$ прост в $K[x]$.

\Proof
Поскольку $f$ неприводим над $K[x]$, то $f$ примитивен, поэтому и неприводим над $F[x]$. Тогда, в силу факториальности $F[x], f$ прост в нем. Если $f|gh$ для некоторых $g, h \in K[x]$, то без ограничения общности $g = gq, q \in F[x]$. Тогда $g = \frac{A}{B}fq_2$, где $q_2 \in K[x]$ -- примитивный многочлен, $\frac{A}{B} \in F$. Тогда аналогично доказанному выделим из многочлена $g$ множитель $c(g)$ и получим, что $f|g$.
\EndProof

Теперь докажем, что $K[x]$ факториально. Теперь это совсем просто. 

В $K[x]$ существует разложение для каждого ненулевого многочлена, и каждый неприводимый над $K[x]$ многочлен прост, поэтому $K[x]$ факториально.
